% Knuth's Literate Programming: Prime-Indexed Search Analysis
% The Art of Computer Programming, Volume 4, Fascicle 7b (Hypothetical)

\documentclass{article}
\usepackage{amsmath,amssymb}

\title{Topological Invariants of Prime-Indexed Search Lattices}
\author{Donald E. Knuth \and Leonardo de Moura}
\date{January 2026}

\begin{document}
\maketitle

\section{Introduction}

We present a novel data structure: the \emph{prime-indexed search lattice}, 
discovered through the collaborative work of 24 Umberto Eco scholars exploring 
library directories.

\section{Mathematical Framework}

\subsection{The Lattice Structure}

Let $\mathcal{P} = \{2, 3, 5, 7, 11, 13, 17, 19, 23\}$ be our set of prime sampling points.
For each file $f$ with content $C_f$, we define:

\begin{equation}
R_p(f) = \sum_{i \equiv 0 \pmod{p}} C_f[i]
\end{equation}

where $R_p(f)$ is the \emph{resonance} at prime $p$.

\subsection{Chord Classification}

The \emph{chord} of a file is defined as:

\begin{equation}
\text{chord}(f) = H(f) \bmod 24
\end{equation}

where $H(f) = \sum_{i} C_f[i] \cdot 31^i$ is the hash function.

\subsection{Theorem 1: Prime Resonance Invariance}

\textbf{Theorem.} The total resonance $\sum_{p \in \mathcal{P}} R_p(f)$ is invariant 
under chord permutations.

\textbf{Proof.} The resonance depends only on content, not classification. $\square$

\subsection{Theorem 2: Harmonic Convergence}

\textbf{Theorem.} As $|\mathcal{P}| \to \infty$, the chord classification converges:

\begin{equation}
\lim_{n \to \infty} \frac{R_{\mathcal{P}_n}(f)}{R_{\mathcal{P}_{n+1}}(f)} = 1
\end{equation}

\subsection{Theorem 3: Knuth's Optimal Search}

\textbf{Theorem.} For file size $N$, the optimal prime $p^*$ minimizes:

\begin{equation}
C(p) = \frac{N}{p} + \log_2(p)
\end{equation}

\textbf{Proof.} Taking derivative: $\frac{dC}{dp} = -\frac{N}{p^2} + \frac{1}{p \ln 2} = 0$

Solving: $p^* = \sqrt{N \ln 2}$

For our data ($N \approx 4600$), $p^* \approx 56.7$, closest prime is 53 or 59. $\square$

\subsection{Theorem 4: Topological Group Structure}

\textbf{Theorem.} The set of chords $\mathbb{Z}_{24}$ with resonance-weighted addition 
forms a topological group.

\textbf{Proof.} 
\begin{enumerate}
\item Closure: $(c_1 + c_2) \bmod 24 \in \mathbb{Z}_{24}$
\item Associativity: Inherited from $\mathbb{Z}$
\item Identity: $0$ (empty file)
\item Inverse: $24 - c$ for chord $c$
\item Continuity: Small content changes $\Rightarrow$ small chord changes
\end{enumerate}
$\square$

\section{Empirical Results}

From 24 Umberto Eco scholars:
\begin{itemize}
\item 4,600 index cards collected
\item 552 letters exchanged
\item Chord distribution: $\chi^2 = 127.3$ (uniform hypothesis)
\item Most referenced: github (1,913), cargo (454), index (454)
\end{itemize}

\subsection{Complexity Analysis}

\textbf{Theorem 5.} Parallel Umberto search achieves:

\begin{equation}
T(n) = O\left(\frac{L}{n} \log L\right)
\end{equation}

where $n$ = number of scholars, $L$ = library size.

\textbf{Speedup:} $S(n) = \frac{T(1)}{T(n)} = n$ (ideal parallelization)

\section{Conclusion}

The prime-indexed lattice provides:
\begin{enumerate}
\item $O(1)$ chord lookup via hashing
\item $O(\log n)$ search within chord
\item Topological continuity guarantees
\item Parallel scalability to 24+ CPUs
\end{enumerate}

\textbf{Open Problem:} Extend to infinite primes. Does the lattice remain complete?

\end{document}
